% GENERATED FILE. Edit canonical lessons, then run npm run generate:books.
% canonical-source-hash: fnv1a64:7fec7fb671965fe4
% canonical-lessons: SA-C495-kisorah, SA-C495-yauvanam, SA-C495-praudhah, SA-C495-varayati, SA-C495-brahmacari

\chapter{Growing Up}
\label{ch:sa-growing-up}
\hlchaptermodality{495}

\begin{chapteropening}
\textbf{By the end of this chapter:} \emph{I can talk about growing up.}

Complete the last lesson of chapter 495: I can talk about growing up.
\end{chapteropening}

\section[\texorpdfstring{\sk{किशोरः} (kiśoraḥ)}{किशोरः (kiśoraḥ)}]{\sk{किशोरः} (kiśoraḥ) — a lad}

\label{lesson:SA-C495-kisorah}



\begin{quote}
Before the new one: say the Sanskrit for permits, then the Sanskrit for childhood.
\end{quote}

\subsection*{You'll want to know: \sk{किशोरः}}
\textbf{\sk{किशोरः}} — \emph{kiśoraḥ} — \textquotedblleft{}a lad\textquotedblright{}.

\begin{quote}
\sk{किशोरः} \sk{क्रीडांगणे} \sk{कन्दुकेन} \sk{क्रीडति।}
\end{quote}

\begin{quote}
\emph{kiśoraḥ krīḍāṅgaṇe kandukena krīḍati.}
\end{quote}

\begin{quote}
The lad plays ball in the playground.
\end{quote}

A masculine noun.

\begin{cousinweb}
\textbf{\sk{किशोरः}} (\emph{kiśoraḥ}) MW's first sense is \textquotedblleft{}a colt\textquotedblright{}, then \textquotedblleft{}a youth, a lad\textquotedblright{}; MW gives no derivation, so its origin is uncertain (MW).
\end{cousinweb}

\begin{grammarlens}[title={What you've built}]
One more word for describing what you lived and felt.
\end{grammarlens}

\subsection*{Your turn}
\begin{itemize}
\raggedright
  \item \emph{Say it:} \emph{kiśoraḥ}
  \item \emph{Say it:} \emph{kiśoraḥ}, once more
  \item \emph{Say it:} \emph{kiśoraḥ krīḍāṅgaṇe kandukena krīḍati.} — The lad plays ball in the playground.
  \item \emph{Recall:} say the Sanskrit for permits, then the Sanskrit for childhood, then say \emph{kiśoraḥ} again
\end{itemize}

\begin{tcolorbox}[breakable,colback=teal!4,colframe=teal!35!black,title={Before you move on}]
What does \sk{किशोरः} mean? (\textquotedblleft{}A lad\textquotedblright{}.) Say it once more.
\end{tcolorbox}
\section[\texorpdfstring{\sk{यौवनम्} (yauvanam)}{यौवनम् (yauvanam)}]{\sk{यौवनम्} (yauvanam) — youth (the time of life)}

\label{lesson:SA-C495-yauvanam}



\begin{quote}
Before the new one: say the Sanskrit for childhood, then the Sanskrit for a lad.
\end{quote}

\subsection*{You'll want to know: \sk{यौवनम्}}
\textbf{\sk{यौवनम्}} — \emph{yauvanam} — \textquotedblleft{}youth (the time of life)\textquotedblright{}.

\begin{quote}
\sk{यौवने} \sk{उत्साहः} \sk{बहु} \sk{भवति।}
\end{quote}

\begin{quote}
\emph{yauvane utsāhaḥ bahu bhavati.}
\end{quote}

\begin{quote}
In youth there is much enthusiasm.
\end{quote}

A neuter noun.

\begin{cousinweb}
\textbf{\sk{यौवनम्}} (\emph{yauvanam}) is built on \emph{yuvan}, \textquotedblleft{}young\textquotedblright{}, the word behind \textbf{\sk{युवकः}}; MW gives \textquotedblleft{}youth, youthfulness, adolescence\textquotedblright{} (MW).
\end{cousinweb}

\begin{grammarlens}[title={What you've built}]
One more word for describing what you lived and felt.
\end{grammarlens}

\subsection*{Your turn}
\begin{itemize}
\raggedright
  \item \emph{Say it:} \emph{yauvanam}
  \item \emph{Say it:} \emph{yauvanam}, once more
  \item \emph{Say it:} \emph{yauvane utsāhaḥ bahu bhavati.} — In youth there is much enthusiasm.
  \item \emph{Recall:} say the Sanskrit for childhood, then the Sanskrit for a lad, then say \emph{yauvanam} again
\end{itemize}

\begin{tcolorbox}[breakable,colback=teal!4,colframe=teal!35!black,title={Before you move on}]
What does \sk{यौवनम्} mean? (\textquotedblleft{}Youth (the time of life)\textquotedblright{}.) Say it once more.
\end{tcolorbox}
\section[\texorpdfstring{\sk{प्रौढः} (prauḍhaḥ)}{प्रौढः (prauḍhaḥ)}]{\sk{प्रौढः} (prauḍhaḥ) — grown up}

\label{lesson:SA-C495-praudhah}



\begin{quote}
Before the new one: say the Sanskrit for a lad, then the Sanskrit for youth (the time of life).
\end{quote}

\subsection*{You'll want to know: \sk{प्रौढः}}
\textbf{\sk{प्रौढः}} — \emph{prauḍhaḥ} — \textquotedblleft{}grown up\textquotedblright{}.

\begin{quote}
\sk{मम} \sk{पुत्री} \sk{इदानीं} \sk{प्रौढा} \sk{अस्ति।}
\end{quote}

\begin{quote}
\emph{mama putrī idānīṃ prauḍhā asti.}
\end{quote}

\begin{quote}
My daughter is grown up now.
\end{quote}

\begin{cousinweb}
\textbf{\sk{प्रौढः}} (\emph{prauḍhaḥ}) is built from \emph{pra} and \emph{ūḍha}, the ppp of the root √vah of \textbf{\sk{वहति}}; MW gives \textquotedblleft{}grown up, full-grown, mature\textquotedblright{} (MW).
\end{cousinweb}

\begin{grammarlens}[title={What you've built}]
One more word for describing what you lived and felt.
\end{grammarlens}

\subsection*{Your turn}
\begin{itemize}
\raggedright
  \item \emph{Say it:} \emph{prauḍhaḥ}
  \item \emph{Say it:} \emph{prauḍhaḥ}, once more
  \item \emph{Say it:} \emph{mama putrī idānīṃ prauḍhā asti.} — My daughter is grown up now.
  \item \emph{Recall:} say the Sanskrit for a lad, then the Sanskrit for youth (the time of life), then say \emph{prauḍhaḥ} again
\end{itemize}

\begin{tcolorbox}[breakable,colback=teal!4,colframe=teal!35!black,title={Before you move on}]
What does \sk{प्रौढः} mean? (\textquotedblleft{}Grown up\textquotedblright{}.) Say it once more.
\end{tcolorbox}
\section[\texorpdfstring{\sk{वारयति} (vārayati)}{वारयति (vārayati)}]{\sk{वारयति} (vārayati) — holds back; forbids}

\label{lesson:SA-C495-varayati}



\begin{quote}
Before the new one: say the Sanskrit for youth (the time of life), then the Sanskrit for grown up.
\end{quote}

\subsection*{You'll want to know: \sk{वारयति}}
\textbf{\sk{वारयति}} — \emph{vārayati} — \textquotedblleft{}holds back; forbids\textquotedblright{}.

\begin{quote}
\sk{वैद्यः} \sk{शर्करां} \sk{वारयति।}
\end{quote}

\begin{quote}
\emph{vaidyaḥ śarkarāṃ vārayati.}
\end{quote}

\begin{quote}
The physician forbids sugar.
\end{quote}

Past (he …): \sk{सः} \sk{शर्करां} \sk{वारितवान्}, \emph{saḥ śarkarāṃ vāritavān}.

\begin{cousinweb}
\textbf{\sk{वारयति}} (\emph{vārayati}) is built on the causative of the root √vṛ, \textquotedblleft{}to cover, to ward off\textquotedblright{}; MW's ppp \emph{vārita} is \textquotedblleft{}warded off, prevented, restrained\textquotedblright{} (MW).
\end{cousinweb}

\begin{grammarlens}[title={What you've built}]
One more word for describing what you lived and felt.
\end{grammarlens}

\subsection*{Your turn}
\begin{itemize}
\raggedright
  \item \emph{Say it:} \emph{vārayati}
  \item \emph{Say it:} \emph{vārayati}, once more
  \item \emph{Say it:} \emph{vaidyaḥ śarkarāṃ vārayati.} — The physician forbids sugar.
  \item \emph{Recall:} say the Sanskrit for youth (the time of life), then the Sanskrit for grown up, then say \emph{vārayati} again
\end{itemize}

\begin{tcolorbox}[breakable,colback=teal!4,colframe=teal!35!black,title={Before you move on}]
What does \sk{वारयति} mean? (\textquotedblleft{}Holds back; forbids\textquotedblright{}.) Say it once more.
\end{tcolorbox}
\section[\texorpdfstring{\sk{ब्रह्मचारी} (brahmacārī)}{ब्रह्मचारी (brahmacārī)}]{\sk{ब्रह्मचारी} (brahmacārī) — a celibate student}

\label{lesson:SA-C495-brahmacari}



\begin{quote}
Before the new one: say the Sanskrit for grown up, then the Sanskrit for holds back.
\end{quote}

\subsection*{You'll want to know: \sk{ब्रह्मचारी}}
\textbf{\sk{ब्रह्मचारी}} — \emph{brahmacārī} — \textquotedblleft{}a celibate student\textquotedblright{}.

\begin{quote}
\sk{ब्रह्मचारी} \sk{गुरोः} \sk{गृहे} \sk{वसति।}
\end{quote}

\begin{quote}
\emph{brahmacārī guroḥ gṛhe vasati.}
\end{quote}

\begin{quote}
The student lives in his teacher's house.
\end{quote}

A masculine noun.

\begin{cousinweb}
\textbf{\sk{ब्रह्मचारी}} (\emph{brahmacārī}) is built from \emph{brahman}, \textquotedblleft{}sacred knowledge\textquotedblright{}, and \emph{cārin}, \textquotedblleft{}practising\textquotedblright{}, from the root √car; MW gives \textquotedblleft{}a young Brāhman who is a student of the Veda, before marriage\textquotedblright{} (MW).
\end{cousinweb}

\begin{grammarlens}[title={What you've built}]
One more word for describing what you lived and felt.
\end{grammarlens}

\subsection*{Your turn}
\begin{itemize}
\raggedright
  \item \emph{Say it:} \emph{brahmacārī}
  \item \emph{Say it:} \emph{brahmacārī}, once more
  \item \emph{Say it:} \emph{brahmacārī guroḥ gṛhe vasati.} — The student lives in his teacher's house.
  \item \emph{Recall:} say the Sanskrit for grown up, then the Sanskrit for holds back, then say \emph{brahmacārī} again
\end{itemize}

\begin{tcolorbox}[breakable,colback=teal!4,colframe=teal!35!black,title={Before you move on}]
What does \sk{ब्रह्मचारी} mean? (\textquotedblleft{}A celibate student\textquotedblright{}.) Say it once more.
\end{tcolorbox}
